\begin{abstract}
Multilevel selection phenomena occur when short-term individual-level reproductive interests conflict with longer-term group-level fitness effects.
Being able to detect and quantify this phenomenon is key in understanding a wide range of important questions around topics ranging from the evolution of multicellularity to the evolution of pathogen virulence.
Within the realm of artificial life research, in particular, the topic of multilevel selection has been highlighted in particular with respect to major transitions in evolution and open-ended evolution.
One approach for detecting multilevel selection dynamics, proposed by \citet{TODO} is in screening for mutations that arise more often in a population than expected by chance (due to transient, individual-level fitness benefits) but are ultimately associated with negative longer-term fitness outcomes (i.e., smaller, shorter-lived descendant clades than expected by chance).
In this work, we use agent-based modeling with known ground truth to assess the efficacy of this approach.
To test these methods under challenging conditions broadly comparable to the original dataset explored by \citet{TODO}, we conduct our experiments using the epidemiological covasim framework, modeling multilevel selection in trade-offs between within-host growth rate and between-host transmissibility.
We also describe an extension to the statistical procedure for screening for clade-level fitness effects, switching from sister-based clade normalization to a more general comparator normalization scheme, necessary to achieve efficacy on in silico data.
afor pathogen phylogenies drawn from host populations of up to 1.2 million agents, we find the method to be sensitive in detecting genome sites under multilevel selection with 30\% effect sizes on fitnes, but do not see sensitivity to smaller 10\% mutation effect sizes.
To test the robustness of this methodology, we conduct additional experiments incorporating extrinsic, time-varying environmental changes and the presence of adaptive turnover in population compositions, and find generally comparable screen performance compared to simpler baseline conditions.
As such, this work represents a promising initial step in ultimately enabing more rigorous quantification of multilevel selection effects that is generalizable across simulation systems, potentiallu also including large-scale decentralized simulations where the necessary data could be collected through sampling approaches.
\end{abstract}
